\chapter{Imaging Workflow}

\section*{Patient focus}
This chapter explains the examination or clinical problem in accessible language. Patients should use it to prepare questions and understand the purpose, experience, benefits, and limitations of imaging; local instructions and the official report take priority.


\label{ch:ch08_workflow}

Every imaging examination follows a defined pathway from the moment a clinician has a question to the moment the result influences patient care. Understanding this workflow helps both patients and clinicians ensure that imaging is safe, efficient, and clinically useful.

---

\section*{1. What to Expect as a Patient --- Step by Step}

If your doctor has ordered an imaging study, here is the typical pathway you will experience:

\begin{enumerate}
    \item \textbf{Referral}: Your doctor writes an order or referral specifying the examination, the clinical question, and relevant history. You may be given a preparation instruction sheet.
    \item \textbf{Booking and scheduling}: The imaging department contacts you (or your referring team) to schedule the examination. Urgent studies are arranged within hours; routine studies may be scheduled days to weeks ahead depending on demand.
    \item \textbf{Preparation}: Some examinations require preparation:
    \begin{itemize}
        \item Fasting for 4--6 hours before CT abdomen with contrast.
        \item Drinking oral contrast solution 1--2 hours before abdominal CT.
        \item Removing all jewellery and metallic objects before MRI.
        \item Kidney blood tests may be needed before some contrast-enhanced studies, depending on the contrast agent, your health and local screening policy.
        \item Arranging a companion for discharge if sedation was used.
    \end{itemize}
    \item \textbf{Arrival and check-in}: At the imaging department, you will confirm your identity (name, date of birth, hospital number). You will be asked about allergies, medications, pregnancy status, and any previous imaging.
    \item \textbf{Changing}: You may be asked to change into a hospital gown and remove all metallic items, including body piercings.
    \item \textbf{The examination itself}: A radiographer (technologist) will position you on the imaging table and explain what will happen. The scan itself may take 5 minutes (X-ray) to 60 minutes (complex MRI). Most scans are painless. You may hear noise during MRI; earplugs or headphones will be provided.
    \item \textbf{Post-examination}: If contrast was used, you will be observed briefly (15--30 minutes) for any adverse reactions. You can usually resume normal activities and diet immediately afterward unless instructed otherwise.
    \item \textbf{Report and results}: The images are interpreted by a radiologist and a report is issued. Your referring doctor will discuss the results with you. Urgent findings are communicated immediately.
\end{enumerate}

\begin{pearlbox}[You can always ask]
If you are anxious, claustrophobic (MRI), pregnant, or unsure about preparation instructions, speak to the imaging team before your appointment. Most examinations can be adapted to your needs.
\end{pearlbox}

---

\section*{2. The Complete Imaging Pathway --- From Clinical Question to Report}

The imaging pathway has 10 stages, each with defined responsibilities and quality checkpoints:

\begin{table}[h]
\centering
\caption{The 10 stages of the imaging pathway}
\begin{tabular}{|l|l|p{5cm}|}
\hline
\textbf{Stage} & \textbf{Responsible} & \textbf{Key activities} \\
\hline
1. Clinical question & Refining clinician & Define the question, review history and prior imaging, determine if imaging is the right next step. \\
\hline
2. Referral / order & Referring clinician & Complete the referral with clinical question, history, relevant findings, and urgency. \\
\hline
3. Protocol selection & Radiologist / protocol team & Select the optimal modality, contrast, phases, and sequences based on the clinical question. \\
\hline
4. Scheduling & Booking team & Match urgency to available slots; ensure preparation instructions are issued. \\
\hline
5. Patient preparation & Radiographer / nurse & Confirm identity, screening (pregnancy, implants, renal function), consent, and preparation. \\
\hline
6. Image acquisition & Radiographer & Positioning, protocol execution, quality checks, dose monitoring. \\
\hline
7. Post-processing & Radiographer / AI tools & 3D reconstruction, MPR, MIP, perfusion maps, AI-assisted detection. \\
\hline
8. Interpretation & Radiologist & Systematic review, comparison with prior studies, AI concurrent reads. \\
\hline
9. Reporting & Radiologist & Structured or narrative report answering the clinical question. \\
\hline
10. Communication & Radiologist / system & Issue report; communicate urgent findings directly to the clinician; follow up on recommendations. \\
\hline
\end{tabular}
\end{table}

---

\section*{3. Writing a Good Referral}
The quality of the imaging study begins with the quality of the referral. A vague referral leads to protocol uncertainty, unnecessary radiation, and inconclusive reports.

\subsection*{Essential Elements of an Imaging Referral}

\begin{table}[h]
\centering
\caption{What to include in an imaging referral}
\begin{tabular}{|l|p{8cm}|}
\hline
\textbf{Element} & \textbf{Why it matters} \\
\hline
Allergies and medications & Affects contrast choice and safety screening. \\
\hline
Brief clinical history & Provides context for interpretation (e.g., ``post-surgical'' vs ``pre-operative''). \\
\hline
Clinical question & Determines the modality, protocol, and contrast. Without it, the radiologist must guess. \\
\hline
Pregnancy status & Affects modality choice and radiation dose considerations. \\
\hline
Prior imaging & Enables comparison and avoids unnecessary repeat imaging. \\
\hline
Relevant examination findings & Guides protocol (e.g., suspected appendicitis triggers appendiceal CT protocol). \\
\hline
Renal function (eGFR) & Screen when indicated by the contrast agent, patient risk and current local policy. \\
\hline
Specific clinical scenario & Refines protocol (e.g., ``pancreatic mass protocol'' vs ``general abdominal CT''). \\
\hline
Urgency & Determines scheduling priority and communication pathway. \\
\hline
\end{tabular}
\end{table}

\subsection*{Sample Referral Template}

\begin{testbox}[IMAGING REFERRAL TEMPLATE]
\textbf{Patient details}: Name, DOB, Hospital ID, Contact phone.\\
\textbf{Clinical question}: ``Is there a pulmonary embolism?''\\
\textbf{Brief history}: 52-year-old male, acute onset dyspnoea and right-sided pleuritic chest pain, D-dimer elevated (4.8).\\
\textbf{Examination findings}: Tachycardic (110 bpm), SpO$_2$ 94\% on room air. Right calf swollen.\\
\textbf{Prior imaging}: CXR today --- clear lung fields.\\
\textbf{Allergies}: Penicillin (rash). No prior contrast reaction.\\
\textbf{Medications}: Warfarin (held 3 days).\\
\textbf{Pregnancy}: N/A.\\
\textbf{Renal function}: eGFR 82 (checked yesterday).\\
\textbf{Urgency}: Urgent --- patient is an inpatient.\\
\textbf{Protocol preference}: CTPA protocol.
\end{testbox}

\begin{pearlbox}[The three-second rule]
If you cannot state your clinical question in one sentence, you are not ready to order the imaging study. Refine your thinking before submitting the referral. ``Abdominal pain, ?cause'' is too vague. ``Right iliac fossa pain with rebound tenderness, ?appendicitis'' is useful.
\end{pearlbox}

---

\section*{4. Protocol Selection}
Once the clinical question is defined, the appropriate protocol is selected. Protocol selection is a partnership between the referring clinician (who knows the clinical context) and the radiologist (who knows the imaging capabilities).

\begin{table}[h]
\centering
\caption{Clinical question to protocol mapping (examples)}
\begin{tabular}{|p{5cm}|l|l|}
\hline
\textbf{Clinical question} & \textbf{Modality} & \textbf{Protocol} \\
\hline
Bone metastases & Nuclear medicine & $^{99m}$Tc MDP bone scan or PET-CT \\
\hline
Knee ligament injury & MRI & Knee MRI (T1, T2, PD fat-sat sequences) \\
\hline
Staging of lung cancer & CT + PET & CT chest with contrast + PET-CT \\
\hline
Suspected appendicitis & CT & CT abdomen/pelvis with IV contrast (portal venous phase) \\
\hline
Suspected cholelithiasis & Ultrasound & Abdominal ultrasound (gallbladder protocol) \\
\hline
Suspected PE & CT & CTPA (arterial phase, timed to pulmonary artery opacification) \\
\hline
Suspected stroke & MRI & Brain MRI (DWI, FLAIR, MRA, SWI) \\
\hline
\end{tabular}
\end{table}

---

\section*{5. The Radiology Report}
The report is the final product of the imaging pathway. It must answer the clinical question, describe findings clearly, and provide actionable recommendations.

\subsection*{Structure of a Radiology Report}
\begin{enumerate}
    \item \textbf{Patient and examination details}: Name, DOB, date of study, type of examination, contrast used.
    \item \textbf{Clinical indication}: The reason for the examination (reproduced from the referral).
    \item \textbf{Technique}: How the study was performed (e.g., ``CT of the chest was performed following IV administration of 80 mL iohexol at 3 mL/s with bolus tracking to the pulmonary artery'').
    \item \textbf{Comparison}: Reference to prior relevant imaging.
    \item \textbf{Findings}: Organ-by-organ or region-by-region description of normal and abnormal findings.
    \item \textbf{Impression}: A concise summary answering the clinical question, with the most important finding first.
    \item \textbf{Recommendations}: Follow-up imaging, further investigation, or clinical correlation as appropriate.
\end{enumerate}

\subsection*{Turnaround Time Expectations}

\begin{table}[h]
\centering
\caption{Typical reporting turnaround times}
\begin{tabular}{|l|l|}
\hline
\textbf{Priority} & \textbf{Target turnaround} \\
\hline
Emergency / on-call & $<$60 minutes (verbal then written) \\
\hline
Routine (outpatient) & 5--14 days \\
\hline
Screening / research & Variable \\
\hline
Urgent (inpatient) & $<$24 hours \\
\hline
\end{tabular}
\end{table}

---

\section*{6. Critical Result Communication}
Some findings require immediate communication to the referring clinician, regardless of the time of day. These are ``critical results'' and must not wait for the formal report.

\begin{testbox}[Examples of critical findings requiring immediate communication]
\begin{itemize}
    \item Tension pneumothorax or large pneumothorax.
    \item New intracranial haemorrhage or mass effect with midline shift.
    \item Pulmonary embolism with right heart strain.
    \item Aortic dissection or ruptured aortic aneurysm.
    \item Free air under the diaphragm (viscus perforation).
    \item Active internal bleeding with haemodynamic instability.
    \item Cauda equina syndrome on MRI.
    \item Ectopic pregnancy on ultrasound.
    \item Spinal cord compression.
    \item Suspected mesenteric ischaemia.
\end{itemize}
\end{testbox}

Critical results must be communicated directly (phone call, in person) and documented with:
\begin{itemize}
    \item The date and time of communication.
    \item The name and designation of the person who communicated the finding.
    \item The name and designation of the person who received the finding.
    \item The content of the communication.
    \item Any agreed action plan.
\end{itemize}

\begin{pearlbox}[The phone call matters more than the report]
A critical finding communicated by phone saves lives. A critical finding left in a report to be read ``later'' can cost lives. If you discover a critical finding, put down the pen and pick up the phone.
\end{pearlbox}

---

\section*{7. Incidental Findings and Follow-Up Recommendations}
Incidental findings are abnormalities detected outside the organ or system being examined. They are common (incidentalomas in up to 40\% of CT studies). Management should be evidence-based:
\begin{itemize}
    \item \textbf{Fleischner Society guidelines} for pulmonary nodules.
    \item \textbf{ACR Incidental Findings Committee} recommendations for adrenal, thyroid, renal, and hepatic incidentalomas.
    \item \textbf{Lung-RADS} for low-dose CT lung cancer screening.
    \item \textbf{BI-RADS}, \textbf{CAD-RADS}, \textbf{PI-RADS}, \textbf{LI-RADS} for organ-specific standardised reporting.
\end{itemize}

Avoid ``cascade investigations''---automated follow-up of every incidental finding without clinical context. Each recommendation should be proportionate and justified.

\clearpage
